\documentclass[11pt]{amsart}
\usepackage[margin=1.1in]{geometry}
\usepackage{amsmath,amssymb,amsthm,mathtools}
\usepackage{enumitem}
\usepackage[hidelinks]{hyperref}

\newtheorem{theorem}{Theorem}[section]
\newtheorem{proposition}[theorem]{Proposition}
\newtheorem{lemma}[theorem]{Lemma}
\newtheorem{corollary}[theorem]{Corollary}
\theoremstyle{definition}
\newtheorem{definition}[theorem]{Definition}
\theoremstyle{remark}
\newtheorem{remark}[theorem]{Remark}

\newcommand{\Wp}{W_{\geq 1}}
\newcommand{\g}{\mathfrak{g}}
\newcommand{\U}{U}
\newcommand{\cJ}{\mathcal{J}}
\newcommand{\cN}{\mathcal{N}}
\newcommand{\cD}{\mathcal{D}}
\newcommand{\gr}{\operatorname{gr}}
\newcommand{\proj}{\operatorname{proj}}
\newcommand{\sym}{\operatorname{sym}}
\newcommand{\ad}{\operatorname{ad}}
\newcommand{\lm}{\operatorname{lm}}
\newcommand{\kk}{\Bbbk}
\newcommand{\eb}{\bar e}

\title[The kernel of $\Psi_1$ on $U(W_{\geq 1})$]{The kernel of the first orbit homomorphism on the\\ positive Witt algebra, via the Conley--Goode method}
\date{}

\begin{document}

\begin{abstract}
Let $\Psi$ be the first orbit homomorphism of Pham--Timmins, restricted to the enveloping algebra of the positive Witt algebra $\Wp$. We prove, with the Conley--Goode criterion, that $\ker\Psi$ is the two-sided ideal generated by the single differentiator $\Omega=\Omega^{(3)}_{5,1}$, recovering a result of Sierra--Walton. The proof is self-contained and uses only elements of positive weight. An exact formula for $\Psi$ shows that $\Psi$ of a two-letter word is quadratic in its index. This gives, in turn: injectivity on normal forms; the quadratic kernel as the span of third differences; and the Conley--Goode condition on quadratic elements, by Lagrange interpolation. Exact bracket formulas for differentiators then show that $\Omega$ generates them all.
\end{abstract}
\maketitle

\section{Statement}

Throughout, $\kk$ is a field of characteristic $0$.

\subsection*{The Lie algebra}
Let $\g=\Wp$ have basis $\{e_i : i\ge1\}$ and bracket $[e_i,e_j]=(j-i)\,e_{i+j}$. Let $\U=\U(\g)$, filtered by order: $\U_m$ is spanned by products of at most $m$ elements of $\g$. By PBW, the nondecreasing monomials $e_{i_1}\cdots e_{i_m}$ form a basis of $\U$, compatible with the filtration, and $S(\g)=\gr\U=\kk[\eb_1,\eb_2,\dots]$. For $z\in\U_m$ write $\gr_mz=z+\U_{m-1}$, and for a subspace $A\subseteq\U$ put $\gr_mA=(A\cap\U_m+\U_{m-1})/\U_{m-1}$. The map $\gr$ is multiplicative, and the symmetrizer $\sym:S(\g)\to\U$ satisfies $\gr_m\sym(x)=x$ for $x\in S^m(\g)$.

Each $\ad e_l$ preserves every $\U_m$, and on $S(\g)$ it induces the derivation $\theta_l$ with $\theta_l(\eb_a)=(a-l)\,\eb_{a+l}$; that is, $\gr_m[e_l,z]=\theta_l(\gr_mz)$ for $z\in\U_m$.

Setting $\deg e_i=i$ (the \emph{weight}) grades $\g$, $\U$ and $S(\g)$, since $[e_i,e_j]\in\kk e_{i+j}$. No element $e_0$ is needed.

\subsection*{The target}
Let $T=A_1[v]=\kk[t,\partial]\otimes\kk[v]$, with $\partial t-t\partial=1$ and $v$ central, graded by weight $\deg t=1$, $\deg\partial=-1$, $\deg v=0$.

\begin{theorem}\label{thm:main}
There is an algebra homomorphism $\Psi:\U\to T$ with $\Psi(e_i)=t^{i+1}\partial+(i+1)t^iv$, and
\[
\ker\Psi=(\Omega),\qquad \Omega=e_5e_1-3e_4e_2+3e_3e_3-e_2e_4=e_1e_5-4e_2e_4+3e_3^2+2e_6 .
\]
\end{theorem}

Here $\Psi$ is the restriction of the first orbit homomorphism $\Psi_1$ of \cite{PT}, and $\Omega=\Omega^{(3)}_{5,1}$ in their notation.

\section{The homomorphism in good coordinates}

Put $X=t\partial+2v\in T$.

\begin{lemma}\label{lem:coords}
\begin{enumerate}[label=\textup{(\roman*)}]
\item $X$ commutes with $v$, and $Xt=t(X+1)$. Hence $g\,t^j=t^j\,g(X+j)$ for every $g\in\kk[X,v]$, where $g(X+j)$ means $g(X+j,v)$.
\item $X$ and $v$ are algebraically independent, and if $t^wg=0$ with $g\in\kk[X,v]$ then $g=0$.
\end{enumerate}
\end{lemma}

\begin{proof}
(i) follows from $(t\partial)t=t(t\partial+1)$. For (ii), note that $\kk[X,v]=\kk[t\partial,v]$, and $(t\partial)^a=t^a\partial^a+(\text{a combination of }t^c\partial^c,\ c<a)$. So the elements $(t\partial)^av^c$ are linearly independent in $T$, which has basis $t^a\partial^bv^c$. Finally, $T$ is a domain and $t^w\neq0$.
\end{proof}

For $i\ge1$ let $\ell_i=X+(i-1)v$ and $Y_i=t^i\ell_i=t^{i+1}\partial+(i+1)t^iv$.

\begin{lemma}\label{lem:formula}
For $i_1,\dots,i_m\ge1$ with $i_1+\cdots+i_m=w$,
\begin{equation}\label{eq:formula}
Y_{i_1}\cdots Y_{i_m}=t^w\prod_{k=1}^m\ell_{i_k}(X+\sigma_k),\qquad \sigma_k=i_{k+1}+\cdots+i_m .
\end{equation}
In particular $[Y_i,Y_j]=(j-i)\,Y_{i+j}$, so $\Psi$ exists, with $\Psi(e_i)=Y_i$.
\end{lemma}

\begin{proof}
Move all powers of $t$ to the left using Lemma~\ref{lem:coords}(i). The $k$-th factor $\ell_{i_k}$ then has $t^{\sigma_k}$ passing it, so it is shifted by $\sigma_k$. For $m=2$ this gives
\[
Y_iY_j-Y_jY_i=t^{i+j}\Bigl[\bigl(X+j+(i-1)v\bigr)\bigl(X+(j-1)v\bigr)-\bigl(X+i+(j-1)v\bigr)\bigl(X+(i-1)v\bigr)\Bigr]=(j-i)\,t^{i+j}\ell_{i+j}.\qedhere
\]
\end{proof}

So for a word $n=e_{i_1}\cdots e_{i_m}$ of weight $w$ we have $\Psi(n)=t^wf_n$, where $f_n=\prod_k\ell_{i_k}(X+\sigma_k)\in\kk[X,v]$.

\begin{corollary}\label{cor:Omega}
$\Psi(\Omega)=0$.
\end{corollary}

\begin{proof}
For $a+b=6$ we have $\Psi(e_ae_b)=t^6F(b)$, where $F(b)=\bigl(X+b+(5-b)v\bigr)\bigl(X+(b-1)v\bigr)$ is a polynomial of degree $2$ in $b$. Hence
\[
\Psi(\Omega)=t^6\bigl(F(1)-3F(2)+3F(3)-F(4)\bigr)=0,
\]
since the third finite difference of a quadratic polynomial vanishes.
\end{proof}

\section{The Conley--Goode criterion}

\begin{theorem}[{\cite[Prop.~3.1]{CG}}]\label{thm:CG}
Let $\Psi:\U(\g)\to T$ be an algebra map and $d\ge1$. Let $J\subseteq\g$, $\cJ\subseteq\U(\g)$ and $I\subseteq\ker\Psi\cap\U_d$ be subspaces such that:
\begin{enumerate}[label=\textup{(\alph*)}]
\item $\Psi$ is injective on $\cJ$;
\item $\U_{d-1}\subseteq\cJ$;
\item $\gr_m\cJ=S^{m-d+1}(J)\,S^{d-1}(\g)$ for all $m\ge d$;
\item $\proj_J(\gr_dI)=S^d(\g/J)$, where $\proj_J:S(\g)\to S(\g/J)$ is induced by $\g\to\g/J$.
\end{enumerate}
Then $\ker\Psi=(I)$.
\end{theorem}

\begin{proof}
Choose $\g=J\oplus C$. Call the number of factors from $C$ in a monomial its \emph{$C$-degree}. The kernel of the algebra map $\proj_J$ is the ideal generated by $J$, so $\ker\proj_J\cap S^d(\g)=J\,S^{d-1}(\g)$.

\emph{Step 1: $S^m(\g)=\gr_m\cJ+S^{m-d}(\g)\,\gr_dI$ for $m\ge d$.} Let $M$ be a monomial of order $m$ and $C$-degree $c$, and induct on $c$. If $c<d$, then $M\in\gr_m\cJ$ by (c). Otherwise write $M=rq$, with $q$ a product of $d$ factors from $C$. By (d), choose $p\in\gr_dI$ with $\proj_J(p)=\proj_J(q)$. Then $q-p\in J\,S^{d-1}(\g)$, so in $M=rp+r(q-p)$ the second term has $C$-degree less than $c$.

\emph{Step 2: $\U=\cJ+(I)$.} We show $\U_m\subseteq\cJ+(I)$ by induction on $m$; the base case $m<d$ is (b). Let $x\in\U_m$ with $m\ge d$. By Step 1, $\gr_mx=\gr_my+\sum_kr_k\gr_dz_k$ with $y\in\cJ$, $r_k\in S^{m-d}(\g)$ and $z_k\in I$. Then $x-y-\sum_k\sym(r_k)z_k\in\U_{m-1}$.

\emph{Step 3.} Clearly $(I)\subseteq\ker\Psi$. Conversely, if $z\in\ker\Psi$, write $z=y+i$ with $y\in\cJ$ and $i\in(I)$. Then $\Psi(y)=0$, so $y=0$ by (a).
\end{proof}

\section{The data, and conditions (b) and (c)}

We take $d=2$ and $J=\operatorname{span}\{e_1,e_2\}$. Let $\cN$ be the set of PBW monomials
\[
e_1^ae_2^b\qquad\text{and}\qquad e_1^ae_2^be_j\quad(j\ge3),
\]
that is, polynomials in $e_1,e_2$ times at most one higher generator, and let $\cJ=\operatorname{span}\cN$. For $I$ we take the whole quadratic part of the kernel,
\[
I=K:=\ker\Psi\cap\U_2 .
\]

Since $\cN$ is part of the PBW basis, $\gr_m\cJ$ is spanned by the symbols of the elements of $\cN$ of length $m$. So
\[
\gr\cJ=\kk[\eb_1,\eb_2]\;\oplus\;\textstyle\bigoplus_{j\ge3}\eb_j\,\kk[\eb_1,\eb_2].
\]
Condition (b) holds because $1,e_1,e_2,e_j\in\cN$. For (c), $S^{m-1}(J)\,S^1(\g)$ is spanned by the monomials $\eb_1^a\eb_2^b\eb_j$ with $a+b=m-1$ and $j\ge1$. Those with $j\le2$ are exactly the $\eb_1^a\eb_2^b$ with $a+b=m$, so this span is $\gr_m\cJ$.

\section{Condition (a)}

\begin{proposition}\label{prop:a}
$\Psi$ is injective on $\cJ$.
\end{proposition}

\begin{proof}
$\Psi$ preserves weight, $T$ is the direct sum of its weight spaces, and $\cN$ consists of weight vectors. By Lemma~\ref{lem:coords}(ii) it therefore suffices to show, for each $w$, that the polynomials $f_n$ with $n\in\cN$ of weight $w$ are linearly independent in $\kk[X,v]$. By \eqref{eq:formula},
\begin{align*}
f_{e_1^ae_2^be_j}&=\prod_{p=0}^{a-1}(X+w-a+p)\cdot\prod_{q=0}^{b-1}(X+j+2q+v)\cdot\bigl(X+(j-1)v\bigr),\\
f_{e_1^ae_2^b}&=\prod_{p=0}^{a-1}(X+w-a+p)\cdot\prod_{q=0}^{b-1}(X+2q+v).
\end{align*}
Order the monomials of $\kk[X,v]$ lexicographically with $v>X$. The leading monomial of $X+c$ is $X$. That of $X+c+v$ is $v$, and so is that of $X+(j-1)v$, because $j\neq1$. Leading monomials multiply, so
\[
\lm f_{e_1^ae_2^be_j}=X^av^{b+1},\qquad \lm f_{e_1^ae_2^b}=X^av^b .
\]
For fixed $w$ these are pairwise distinct. A monomial $X^av^k$ is the leading monomial of at most one mixed normal form, namely the one with $b=k-1$ and $j=w-a-2b$. It is also the leading monomial of at most one pure normal form, namely the one with $b=k$. It cannot come from both, because then $j=w-a-2(k-1)=2$. Polynomials with distinct leading monomials are linearly independent.
\end{proof}

\section{The quadratic kernel}

For $k\ge4$ and $s\ge1$ put $c_i=(-1)^i\binom3i$ (with $c_i=0$ for $i\notin\{0,1,2,3\}$) and
\[
D_{k,s}=\sum_{i}c_i\,e_{k-i}e_{s+i}=e_ke_s-3e_{k-1}e_{s+1}+3e_{k-2}e_{s+2}-e_{k-3}e_{s+3},
\]
a combination of words of weight $k+s$. So $D_{k,s}=\Omega^{(3)}_{k,s}$, and $\Omega=D_{5,1}$.

\begin{lemma}\label{lem:F}
Fix $w\ge3$. For $1\le c\le w-1$ we have $\Psi(e_{w-c}e_c)=t^wF(c)$, where
\[
F(c)=F_0+c\,F_1+c^2F_2,\qquad
F_0=\bigl(X+(w-1)v\bigr)(X-v),\quad F_1=X-v+wv^2,\quad F_2=v-v^2,
\]
and $F_0,F_1,F_2$ are linearly independent in $\kk[X,v]$.
\end{lemma}

\begin{proof}
By \eqref{eq:formula}, $F(c)=\bigl(X+c+(w-c-1)v\bigr)\bigl(X+(c-1)v\bigr)$. Expanding in powers of $c$ gives the stated $F_0,F_1,F_2$. Only $F_0$ contains $X^2$, and only $F_1$ contains $X$ without $v$, while $F_2\neq0$. So they are independent.
\end{proof}

\begin{proposition}\label{prop:K}
Let $K=\ker\Psi\cap\U_2$. For every $w$, the weight-$w$ part $K_w$ is spanned by the $D_{w-s,s}$ with $1\le s\le w-4$. In particular every $D_{k,s}$ lies in $\ker\Psi$.
\end{proposition}

\begin{proof}
For $w\le2$ both sides are $0$, since $\Psi$ is injective on $\U_1\subseteq\cJ$. Let $w\ge3$. Every element of $(\U_2)_w$ is a combination of words $e_{w-c}e_c$, because $e_w=(e_{w-1}e_1-e_1e_{w-1})/(2-w)$. Consider the linear map
\[
\lambda=(\lambda_1,\dots,\lambda_{w-1})\in\kk^{w-1}\ \longmapsto\ z_\lambda=\sum_c\lambda_c\,e_{w-c}e_c\in(\U_2)_w ,
\]
which is onto. By Lemma~\ref{lem:F},
\[
\Psi(z_\lambda)=t^w\Bigl(F_0\sum_c\lambda_c+F_1\sum_c\lambda_cc+F_2\sum_c\lambda_cc^2\Bigr).
\]
So $z_\lambda\in\ker\Psi$ if and only if $\lambda$ lies in
\[
V=\Bigl\{\lambda:\ \textstyle\sum_c\lambda_c\,c^r=0\ \text{ for } r=0,1,2\Bigr\}.
\]
The three conditions are independent (Vandermonde, as $w-1\ge3$), so $\dim V=w-4$. For $1\le s\le w-4$, let $\delta^{(s)}$ have entries $c_0,c_1,c_2,c_3$ in positions $s,\dots,s+3$ and $0$ elsewhere. Then $\delta^{(s)}\in V$, because a third finite difference kills polynomials of degree $\le2$. The $\delta^{(s)}$ are independent, since $\delta^{(s)}$ is the first one with a nonzero entry in position $s$. Hence $V=\operatorname{span}\{\delta^{(s)}\}$. Finally $z_{\delta^{(s)}}=D_{w-s,s}$, so $K_w=z(V)=\operatorname{span}\{D_{w-s,s}\}$.
\end{proof}

\section{Condition (d)}

\begin{proposition}\label{prop:d}
$\proj_J(\gr_2K)=S^2(\g/J)$.
\end{proposition}

\begin{proof}
$S^2(\g/J)$ is spanned by the classes of the products $\eb_a\eb_b$ with $3\le a\le b$. Fix such $a,b$ and put $w=a+b\ge6$. By Lemma~\ref{lem:F}, $F$ is a polynomial of degree $2$, and the nodes $1,2,w-1$ are distinct. So Lagrange interpolation gives $F(b)=\sum_{c\in\{1,2,w-1\}}L_c(b)F(c)$, where $L_c$ are the Lagrange basis polynomials. Hence
\[
z=e_ae_b-L_1(b)\,e_{w-1}e_1-L_2(b)\,e_{w-2}e_2-L_{w-1}(b)\,e_1e_{w-1}\ \in K .
\]
Each subtracted word contains $e_1$ or $e_2$, so $\gr_2z\equiv\eb_a\eb_b$ modulo $J\,S^1(\g)$, and $\proj_J(\gr_2z)$ is the class of $\eb_a\eb_b$.
\end{proof}

Conditions (a)--(d) hold by Proposition~\ref{prop:a}, Section~4 and Proposition~\ref{prop:d}, with $I=K\subseteq\ker\Psi\cap\U_2$. Theorem~\ref{thm:CG} therefore gives
\begin{equation}\label{eq:quadgen}
\ker\Psi=(K).
\end{equation}

\section{The differentiators are generated by \texorpdfstring{$\Omega$}{Omega}}

\begin{lemma}\label{lem:brackets}
For $k\ge4$ and $s\ge1$:
\begin{enumerate}[label=\textup{(\roman*)}]
\item $D_{k,s}=-D_{s+3,k-3}$. In particular $D_{4,1}=D_{5,2}=0$ and $D_{4,2}=-D_{5,1}$.
\item $[e_1,D_{k,s}]=(k-1)\,D_{k+1,s}+(s+2)\,D_{k,s+1}$.
\item $[e_2,D_{k,s}]=(k-2)\,D_{k+2,s}+3\,D_{k+1,s+1}+(s+1)\,D_{k,s+2}$.
\end{enumerate}
\end{lemma}

\begin{proof}
(i) Reindexing $i\mapsto3-i$ gives $D_{s+3,k-3}=-\sum_ic_i\,e_{s+i}e_{k-i}$. Hence
\[
D_{k,s}+D_{s+3,k-3}=\sum_ic_i\,[e_{k-i},e_{s+i}]=\Bigl(\sum_ic_i\,(s-k+2i)\Bigr)e_{k+s}=0,
\]
since $\sum_ic_i=\sum_ic_i\,i=0$.

(ii) and (iii). For words, $[e_l,e_ae_b]=(a-l)\,e_{a+l}e_b+(b-l)\,e_ae_{b+l}$ exactly. So both sides of (ii) are combinations of the words $e_{k+1-r}e_{s+r}$, and comparing coefficients, (ii) amounts to
\[
c_r(k-r-1)+c_{r-1}(s+r-2)=(k-1)c_r+(s+2)c_{r-1},\quad\text{i.e.}\quad r\,c_r=(r-4)\,c_{r-1}\qquad(\text{all }r).
\]
This holds because $\binom3r/\binom3{r-1}=(4-r)/r$. Similarly, comparing coefficients of $e_{k+2-r}e_{s+r}$, (iii) amounts to
\[
r\,c_r+3\,c_{r-1}=(r-5)\,c_{r-2}\qquad(\text{all }r),
\]
which one checks for $r=0,\dots,5$, using $(c_0,\dots,c_3)=(1,-3,3,-1)$.
\end{proof}

\begin{proposition}\label{prop:gen}
Let $M$ be the $\ad(\g)$-submodule of $\U$ generated by $\Omega$. Then $D_{k,s}\in M$ for all $k\ge4$ and $s\ge1$.
\end{proposition}

\begin{proof}
We induct on the weight $w=k+s$. The weight-$w$ differentiators are $D_{w-s,s}$ with $1\le s\le w-4$.

\emph{Weights $\le7$.} By Lemma~\ref{lem:brackets}(i), $D_{4,1}=0$, $D_{4,2}=-\Omega$ and $D_{5,2}=0$. By (ii), $[e_1,\Omega]=4D_{6,1}+3D_{5,2}=4D_{6,1}$, so $D_{6,1}\in M$, and then $D_{4,3}=-D_{6,1}\in M$.

\emph{From weights $w-1,w$ to $w+1$, for $w\ge7$.} Put $X_0=D_{w,1}$. For $1\le s\le w-4$, (ii) gives
\[
[e_1,D_{w-s,s}]=(w-s-1)\,D_{w+1-s,s}+(s+2)\,D_{w-s,s+1}\in M .
\]
By induction on $s$, $D_{w+1-s,s}\equiv\mu_s\,X_0\pmod M$ for $1\le s\le w-3$, where $\mu_1=1$ and $\mu_{s+1}=-\frac{w-s-1}{s+2}\,\mu_s$. In particular $\mu_2=-\frac{w-2}{3}$ and $\mu_3=\frac{(w-2)(w-3)}{12}$. Now (iii) applied to $D_{w-2,1}\in M$ gives
\begin{align*}
M\ni[e_2,D_{w-2,1}]&=(w-4)\,X_0+3\,D_{w-1,2}+2\,D_{w-2,3}\\
&\equiv\bigl(w-4+3\mu_2+2\mu_3\bigr)X_0=\frac{(w-6)(w+1)}{6}\,X_0 .
\end{align*}
The coefficient is nonzero for $w\ge7$, so $X_0\in M$, and therefore every $D_{w+1-s,s}\in M$.
\end{proof}

\section{Proof of Theorem~\ref{thm:main}}

By \eqref{eq:quadgen}, Proposition~\ref{prop:K} and Proposition~\ref{prop:gen},
\[
\ker\Psi=(K)=\bigl(D_{k,s}:k\ge4,\ s\ge1\bigr)\subseteq(M)\subseteq(\Omega)\subseteq\ker\Psi ,
\]
using $[e_l,z]=e_lz-ze_l$ for $M\subseteq(\Omega)$, and Corollary~\ref{cor:Omega} for the last inclusion. Hence $\ker\Psi=(\Omega)$.\qed

\begin{remark}\leavevmode
\begin{enumerate}[label=(\arabic*)]
\item The interpolation argument of Proposition~\ref{prop:d} also works in Pham--Timmins' setting $\U(W_{\ge-1})$, with $J=\operatorname{span}\{e_{-1},\dots,e_{n-1}\}$ and $I=\ker\Psi_n\cap\U_2$. There $\Psi_n(e_{w-c}e_c)$ is a polynomial in $c$ of degree at most $2n$, and the $2n+2$ words with a letter in $J$ supply enough interpolation nodes. (We checked the resulting statement by machine for $n\le3$ and weights $\le14$.)
\item Lemma~\ref{lem:brackets}, Proposition~\ref{prop:K} (in weights $\le14$) and the elements $z$ of Proposition~\ref{prop:d} were also confirmed by machine.
\end{enumerate}
\end{remark}

\begin{thebibliography}{9}
\bibitem{CG} C.~H.~Conley and W.~Goode, \emph{An approach to annihilators in the context of vector field Lie algebras}, Expo.\ Math.\ \textbf{43} (2025), no.~2, Paper No.~125600.
\bibitem{PT} T.~A.~Pham and J.~Timmins, \emph{The kernel and image of orbit homomorphisms for the Witt algebra}, arXiv:2510.00756.
\bibitem{SW} S.~J.~Sierra and C.~Walton, \emph{Maps from the enveloping algebra of the positive Witt algebra to regular algebras}, Pacific J.\ Math.\ \textbf{284} (2016), no.~2, 475--509.
\end{thebibliography}

\end{document}
